\documentclass[10pt,a4paper]{amsart}
\usepackage{amsmath,amssymb,amsthm,graphicx}
\numberwithin{equation}{section}

\newtheorem{Theorem}{Theorem}[section]
\newtheorem*{Theorem*}{Theorem}
\newtheorem{Corollary}[Theorem]{Corollary}
\newtheorem{Lemma}[Theorem]{Lemma}
\newtheorem{Proposition}[Theorem]{Proposition}
\newtheorem{Conjecture}[Theorem]{Conjecture}
 { \theoremstyle{definition}
\newtheorem{Definition}[Theorem]{Definition}
\newtheorem{Note}[Theorem]{Note}
\newtheorem{Example}[Theorem]{Example}
\newtheorem{Remark}[Theorem]{Remark} }

\newcommand{\leftexp}[2]{{\vphantom{#2}}^{#1}{#2}}

\usepackage{amscd,amsxtra}
\usepackage{mathtools}
\usepackage{tensor}
\usepackage{faktor}
\usepackage{dsfont}
\usepackage[all]{xy}
\usepackage{tikz, tikz-cd}
\usetikzlibrary{arrows, matrix, intersections, math, calc, decorations.pathmorphing, decorations.markings, positioning}
\usepackage{enumitem}
\newcommand{\CC}{\mathbb{C}}
\newcommand{\ZZ}{\mathbb{Z}}
\newcommand{\PP}{\mathbb{P}}
\newcommand{\RR}{\mathbb{R}}
\newcommand{\QQ}{\mathbb{Q}}
\newcommand{\ii}{\mathbf{i}}
\newcommand{\one}{\mathbf{1}}
\newcommand{\op}{\operatorname}
\newcommand{\A}{\mathcal{A}}
\renewcommand{\H}{\mathcal{H}}
\newcommand{\T}{\mathcal{T}}
\newcommand{\C}{\mathcal{C}}
\renewcommand{\O}{\mathcal{O}}
\newcommand{\D}{\mathcal{D}}
\newcommand{\F}{\mathcal{F}}
\newcommand{\M}{\mathcal{M}}
\newcommand{\E}{\mathcal{E}}
\newcommand{\I}{\mathcal{I}}
\usepackage{cleveref}
\crefname{originalproposition}{original proposition}{original propositions}
\Crefname{originalproposition}{Original proposition}{Original propositions}
\makeatletter
\newlabel{prop:tw=proj}{{internal}{source}{}{sourcepointer-1}{}}
\newlabel{prop:tw=proj@cref}{{[originalproposition][0][]internal}{[1][0][]source}{}{}{}}
\makeatother
\title{Reflection classes of the exceptional-divisor period are proportional to the Iritani images of the exceptional line-bundle differences}
\author{Milanov, Todor, Xia, Xiaokun}
\date{doi:10.3842/SIGMA.2024.029; CC BY 4.0}
\begin{document}
\maketitle
\setcounter{section}{5}
\setcounter{Theorem}{6}
\noindent\textit{Editorial note.} All delivered mathematical text is the author's own, copied from the pinned SIGMA article source; no statement or proof step was rewritten, reordered or removed. The counted statement is the labeled corollary environment at source lines 3229-3241 and the counted proof is the proof environment at source lines 3243-3282, which immediately follows it (source line 3242 between them is blank) and is its own proof.\par\medskip
\noindent\textit{Editorial note.} The proof invokes Proposition prop:tw=proj of the same chapter once; it is a navigation stub with its role recorded in the reason, and it is not a bibliography entry. The reflection loops gamma\_k(q) referenced in the statement are defined by the author in the immediately preceding discussion of the same section, which the proof text takes as given.\par\medskip
\noindent\textit{Editorial note.} The counted statement and proof cite no bibliography entry, so this card needs no bibliography.\par\medskip
\noindent\textit{Editorial note.} The pinned source uses the publisher class sigma, which is not part of the assembly toolchain; the retained preamble is a mechanical, line-for-line copy of the author's package, macro and theorem declaration block (pinned source lines 9-52 minus the publisher-class and page-setup lines). Nothing was added, removed or redefined, and the excerpt serves only as preamble, never as a span coordinate.\par\medskip
\noindent\textit{Allowed standard prerequisites (editorial).}\begin{itemize}
\item The period comparison between the twisted exceptional-divisor periods and the projective periods, Proposition prop:tw=proj of this chapter. invoked at proof line 3244. Transports the reflection condition on the twisted period to the proportionality of sigma(beta) with the projective reflection vector.
\item The Iritani integral-structure map, whose value on a line-bundle difference is expressed through the Gamma class and the Chern character. used at proof lines 3261-3281. Provides the target expression for the proportionality check: the Gamma class of the blowup factors through the Gamma classes of the base and the exceptional directions, and the Chern character of O((-k+1)E)-O(-kE) contributes the exponential factor.
\end{itemize}\medskip
\ifdefined\hypertarget\hypertarget{sourcepointer-1}{}\fi
\noindent\textit{Source reference.} Original proposition internal: sources/sigma-2024-029-reflection-vectors-and-quantum-cohomology-of/pinned-original-source.tex, source lines 3096--3106. The reference is invoked at proof line 3244 to transport the reflection condition to the projective side; it is the author's own preceding result in the same chapter, delivered as its own card, and the proof text here is preserved verbatim as the author wrote it.\par\medskip
\begin{Corollary}\label{cor:tw_refl}
If $\beta\in \widetilde{H}(E)$ is such that the analytic continuation
of \smash{$\leftexp{tw}{I}^{(-m)}_\beta(Q,\lambda)$} along $\gamma_k(q)$ is
\smash{$\leftexp{tw}{I}^{(-m)}_{-\beta}(Q,\lambda)$}, then $\beta$ must be
proportional to
\[
\Psi(\O_E(-k+1)) = (2\pi)^{\frac{1-n}{2}}
\Gamma(\operatorname{Bl}(X)) Q^{-e(n-1)}
(2\pi\ii)^{\operatorname{deg}}
\operatorname{ch}(\O_E(-k+1)),
\]
where $\O_E(-k+1):=\O(-(k-1) E)-\O(-kE)$.
\end{Corollary}

\begin{proof}
According to the above discussion and Proposition \ref{prop:tw=proj},
under the isomorphism $\widetilde{H}(E)\cong H\big(\PP^{n-2}\big)$,
$\sigma(\beta)$ must be proportional to $\Psi_q(\O(k))$, that is,
$\beta$ is proportional to
\[
e^{-\pi\ii \theta} \
(2\pi)^{\frac{3-n}{2}} \Gamma(1+p)^{n-1} q^{-p} e^{2\pi k p}=
(2\pi)^{(3-n)/2} \ii^{2-n}
\Gamma(1-p)^{n-1} q^p e^{-2\pi\ii k p}.
\]
Note that $q^p = e^{\pi\ii p} Q^{-p(n-1)}$. Therefore, under the
isomorphism $\widetilde{H}(E)\cong H\big(\PP^{n-2}\big)$, the above
expression becomes
\[
(2\pi)^{(3-n)/2} \ii^{2-n}
\Gamma(1-e)^{n-1} Q^{-e(n-1)} e^{(-2k+1)\pi\ii e} e.
\]
We have to check that the above expression is proportional to the
image of the Iritani map for~$\operatorname{Bl}(X)$ of the exceptional object
$\O((-k+1)E)-\O(-kE).$ We have
\begin{gather*}
\Gamma(\operatorname{Bl}(X))= \Gamma(X) \Gamma(1-e)^n \Gamma(1+e),\\
\Gamma(1-e)\Gamma(1+e) = \frac{2\pi\ii e}{e^{\pi\ii e}-e^{-\pi\ii
 e}} =
\frac{2\pi\ii e}{e^{2\pi\ii e}-1} e^{\pi\ii e},
\end{gather*}
and
\[
(2\pi\ii)^{\operatorname{deg}}
\operatorname{ch}(\O((-k+1)E)-\O(-kE))=
e^{-2\pi\ii k e }\big(e^{2\pi\ii e}-1\big).
\]
Since $\Gamma(X)\cup e= e$, the image of the Iritani map becomes
\[
(2\pi)^{(3-n)/2} \ii \Gamma(1-e)^{n-1} Q^{-e(n-1)}
e^{(-2k+1)\pi\ii e} e.
\]
The claim of the lemma follows.
\end{proof}

\end{document}
